\subsection{Iterable vs. Iterator}

The words \emph{iterable} and \emph{iterator} are closely related, but they do not name the same thing. Confusing them makes Python loops look more mysterious than they really are.

An iterable is an object that can provide a traversal object. An iterator is the traversal object itself. The iterable is the source. The iterator is the moving cursor over that source.

\begin{quote}
An iterable can be asked to begin traversal. An iterator is the object that performs the traversal step by step.
\end{quote}

This distinction explains why a list, tuple, string, and \texttt{range} can be traversed many times, while a specific iterator is consumed as it advances.

\subsubsection{Iterable Objects: Objects That Can Produce an Iterator via \texttt{iter()}}

An iterable object is an object that can be passed to \texttt{iter()}. If the call succeeds, Python receives an iterator that can produce values one at a time.

Examples of common iterable objects include strings, lists, tuples, dictionaries, sets, and \texttt{range} objects. In this chapter, the most important examples are strings, lists, tuples, and ranges.

\begin{verbatim}
text = "Ni!"
items = ["soup", "pretzels", "muffin"]
coords = (10, 20, 30)
steps = range(1, 4)

print(f"type(text): {type(text)}")
print(f"type(items): {type(items)}")
print(f"type(coords): {type(coords)}")
print(f"type(steps): {type(steps)}")

print(f"has __iter__ text: {'__iter__' in dir(text)}")
print(f"has __iter__ items: {'__iter__' in dir(items)}")
print(f"has __iter__ coords: {'__iter__' in dir(coords)}")
print(f"has __iter__ steps: {'__iter__' in dir(steps)}")
\end{verbatim}

Possible output:

\begin{verbatim}
type(text): <class 'str'>
type(items): <class 'list'>
type(coords): <class 'tuple'>
type(steps): <class 'range'>
has __iter__ text: True
has __iter__ items: True
has __iter__ coords: True
has __iter__ steps: True
\end{verbatim}

The presence of \texttt{\_\_iter\_\_} is the important signal here. It means the object can produce an iterator.

\begin{verbatim}
items = ["soup", "pretzels", "muffin"]

it = iter(items)

print(f"items: {items}")
print(f"type(items): {type(items)}")
print(f"type(it): {type(it)}")
print(f"items is it: {items is it}")
\end{verbatim}

Possible output:

\begin{verbatim}
items: ['soup', 'pretzels', 'muffin']
type(items): <class 'list'>
type(it): <class 'list_iterator'>
items is it: False
\end{verbatim}

The list and the iterator are different objects. The list stores the elements. The iterator stores the traversal position.

The same iterable can produce more than one iterator:

\begin{verbatim}
items = ["soup", "pretzels", "muffin"]

it1 = iter(items)
it2 = iter(items)

print(f"type(it1): {type(it1)}")
print(f"type(it2): {type(it2)}")
print(f"it1 is it2: {it1 is it2}")

print(f"next(it1): {next(it1)}")
print(f"next(it1): {next(it1)}")
print(f"next(it2): {next(it2)}")
\end{verbatim}

Possible output:

\begin{verbatim}
type(it1): <class 'list_iterator'>
type(it2): <class 'list_iterator'>
it1 is it2: False
next(it1): soup
next(it1): pretzels
next(it2): soup
\end{verbatim}

Advancing \texttt{it1} does not advance \texttt{it2}. They are separate traversal objects over the same list.

A \texttt{range} behaves the same way:

\begin{verbatim}
steps = range(1, 4)

it1 = iter(steps)
it2 = iter(steps)

print(f"type(steps): {type(steps)}")
print(f"type(it1): {type(it1)}")
print(f"steps is it1: {steps is it1}")

print(f"next(it1): {next(it1)}")
print(f"next(it2): {next(it2)}")
\end{verbatim}

Possible output:

\begin{verbatim}
type(steps): <class 'range'>
type(it1): <class 'range_iterator'>
steps is it1: False
next(it1): 1
next(it2): 1
\end{verbatim}

The \texttt{range} object remains a compact arithmetic descriptor. Each call to \texttt{iter(steps)} creates a fresh iterator over that descriptor.

For an iterable, some important methods are:

\begin{itemize}
    \item \texttt{\_\_iter\_\_()}: produces an iterator;
    \item \texttt{\_\_len\_\_()}: often gives the number of elements or described values;
    \item \texttt{\_\_getitem\_\_()}: often supports indexed access, depending on the object type;
    \item \texttt{\_\_contains\_\_()}: often supports membership tests with \texttt{in}.
\end{itemize}

Example with a tuple:

\begin{verbatim}
coords = (10, 20, 30)

print(f"type(coords): {type(coords)}")
print(f"dir(coords) sample: {dir(coords)[:12]}")

print(f"len(coords): {len(coords)}")
print(f"coords[1]: {coords[1]}")
print(f"20 in coords: {20 in coords}")

print(f"coords.__iter__(): {coords.__iter__()}")
print(f"coords.__len__(): {coords.__len__()}")
print(f"coords.__getitem__(1): {coords.__getitem__(1)}")
print(f"coords.__contains__(20): {coords.__contains__(20)}")
\end{verbatim}

Possible output:

\begin{verbatim}
type(coords): <class 'tuple'>
dir(coords) sample: ['__add__', '__class__', '__class_getitem__',
'__contains__', '__delattr__', '__dir__', '__doc__', '__eq__',
'__format__', '__ge__', '__getattribute__', '__getitem__']
len(coords): 3
coords[1]: 20
20 in coords: True
coords.__iter__(): <tuple_iterator object at 0x...>
coords.__len__(): 3
coords.__getitem__(1): 20
coords.__contains__(20): True
\end{verbatim}

The exact memory address printed for the iterator will differ between executions. The important fact is that \texttt{coords.\_\_iter\_\_()} returns a \texttt{tuple\_iterator} object.

\subsubsection{Iterator Objects: Objects That Return Successive Values via \texttt{next()}}

An iterator is an object that can return the next value in a traversal. It does this through \texttt{\_\_next\_\_()}, which is called by the built-in \texttt{next()} function.

\begin{verbatim}
items = ["soup", "pretzels", "muffin"]
it = iter(items)

print(f"type(it): {type(it)}")
print(f"dir(it) sample: {dir(it)[:12]}")
print(f"has __iter__: {'__iter__' in dir(it)}")
print(f"has __next__: {'__next__' in dir(it)}")

print(f"next(it): {next(it)}")
print(f"next(it): {next(it)}")
print(f"next(it): {next(it)}")
\end{verbatim}

Possible output:

\begin{verbatim}
type(it): <class 'list_iterator'>
dir(it) sample: ['__class__', '__delattr__', '__dir__', '__doc__',
'__eq__', '__format__', '__ge__', '__getattribute__', '__getstate__',
'__gt__', '__hash__', '__init__']
has __iter__: True
has __next__: True
next(it): soup
next(it): pretzels
next(it): muffin
\end{verbatim}

Each call to \texttt{next(it)} advances the iterator. The iterator remembers where it is. That position is internal state.

The special method call is visible directly:

\begin{verbatim}
items = ["soup", "pretzels"]
it = iter(items)

print(f"it.__next__(): {it.__next__()}")
print(f"next(it): {next(it)}")
\end{verbatim}

Possible output:

\begin{verbatim}
it.__next__(): soup
next(it): pretzels
\end{verbatim}

The built-in \texttt{next(it)} and the direct call \texttt{it.\_\_next\_\_()} advance the same iterator. They are not separate traversal channels.

An iterator also has \texttt{\_\_iter\_\_()}, but for an iterator this method usually returns the iterator itself:

\begin{verbatim}
items = ["soup", "pretzels"]
it = iter(items)

same_it = iter(it)

print(f"type(it): {type(it)}")
print(f"type(same_it): {type(same_it)}")
print(f"it is same_it: {it is same_it}")

print(f"next(it): {next(it)}")
print(f"next(same_it): {next(same_it)}")
\end{verbatim}

Possible output:

\begin{verbatim}
type(it): <class 'list_iterator'>
type(same_it): <class 'list_iterator'>
it is same_it: True
next(it): soup
next(same_it): pretzels
\end{verbatim}

This is a crucial distinction:

\begin{itemize}
    \item Calling \texttt{iter()} on a list creates a new iterator.
    \item Calling \texttt{iter()} on a list iterator returns the same iterator.
\end{itemize}

That is why iterators are not automatically reusable.

The same idea appears with \texttt{range} iterators:

\begin{verbatim}
steps = range(1, 4)
it = iter(steps)

print(f"type(steps): {type(steps)}")
print(f"type(it): {type(it)}")

print(f"iter(steps) is iter(steps): {iter(steps) is iter(steps)}")
print(f"iter(it) is it: {iter(it) is it}")
\end{verbatim}

Possible output:

\begin{verbatim}
type(steps): <class 'range'>
type(it): <class 'range_iterator'>
iter(steps) is iter(steps): False
iter(it) is it: True
\end{verbatim}

The \texttt{range} object can produce fresh iterators. The \texttt{range\_iterator} is already the moving traversal object.

This explains the exact mechanism behind a \texttt{for} loop:

\begin{verbatim}
for value in iterable:
    body
\end{verbatim}

is conceptually:

\begin{verbatim}
_iterator = iter(iterable)

while True:
    try:
        value = next(_iterator)
    except StopIteration:
        break

    body
\end{verbatim}

The \texttt{for} loop does not need a length. It needs an iterator that either produces a next value or raises \texttt{StopIteration}.

\subsubsection{Iterator Exhaustion: Why Some Traversal Objects Cannot Be Reused After Completion}

Iterator exhaustion means that an iterator has reached the end of its traversal stream. Once exhausted, further calls to \texttt{next()} keep raising \texttt{StopIteration}.

\begin{verbatim}
items = ["soup", "pretzels"]
it = iter(items)

print(f"next(it): {next(it)}")
print(f"next(it): {next(it)}")

try:
    print(next(it))
except StopIteration as err:
    print(f"type(err): {type(err)}")
    print("iterator is exhausted")
\end{verbatim}

Possible output:

\begin{verbatim}
next(it): soup
next(it): pretzels
type(err): <class 'StopIteration'>
iterator is exhausted
\end{verbatim}

The iterator does not automatically restart. It remains at the exhausted position.

\begin{verbatim}
items = ["soup", "pretzels"]
it = iter(items)

for item in it:
    print(f"first traversal: {item}")

for item in it:
    print(f"second traversal: {item}")

print("done")
\end{verbatim}

Possible output:

\begin{verbatim}
first traversal: soup
first traversal: pretzels
done
\end{verbatim}

The second loop prints nothing because it receives the same already-exhausted iterator.

By contrast, looping over the original list twice works:

\begin{verbatim}
items = ["soup", "pretzels"]

for item in items:
    print(f"first traversal: {item}")

for item in items:
    print(f"second traversal: {item}")
\end{verbatim}

Possible output:

\begin{verbatim}
first traversal: soup
first traversal: pretzels
second traversal: soup
second traversal: pretzels
\end{verbatim}

The list is reusable because each \texttt{for} loop calls \texttt{iter(items)} and receives a fresh iterator.

The same difference can be shown with \texttt{range}:

\begin{verbatim}
steps = range(1, 4)

for n in steps:
    print(f"first range traversal: {n}")

for n in steps:
    print(f"second range traversal: {n}")
\end{verbatim}

Possible output:

\begin{verbatim}
first range traversal: 1
first range traversal: 2
first range traversal: 3
second range traversal: 1
second range traversal: 2
second range traversal: 3
\end{verbatim}

But if the iterator itself is reused:

\begin{verbatim}
steps = range(1, 4)
it = iter(steps)

for n in it:
    print(f"first iterator traversal: {n}")

for n in it:
    print(f"second iterator traversal: {n}")

print("finished")
\end{verbatim}

Possible output:

\begin{verbatim}
first iterator traversal: 1
first iterator traversal: 2
first iterator traversal: 3
finished
\end{verbatim}

Again, the second traversal prints nothing because the same iterator has already been consumed.

This matters when assigning iterators to names:

\begin{verbatim}
items = ["Jerry", "Elaine", "Kramer"]

it = iter(items)

first = list(it)
second = list(it)

print(f"first: {first}")
print(f"second: {second}")
\end{verbatim}

Possible output:

\begin{verbatim}
first: ['Jerry', 'Elaine', 'Kramer']
second: []
\end{verbatim}

The first \texttt{list(it)} consumes the iterator. The second \texttt{list(it)} receives no remaining values.

The type of the materialized result is separate from the iterator:

\begin{verbatim}
items = ["Jerry", "Elaine", "Kramer"]
it = iter(items)
first = list(it)

print(f"type(items): {type(items)}")
print(f"type(it): {type(it)}")
print(f"type(first): {type(first)}")

print(f"first: {first}")
print(f"has append on first: {'append' in dir(first)}")
print(f"has __next__ on first: {'__next__' in dir(first)}")
\end{verbatim}

Possible output:

\begin{verbatim}
type(items): <class 'list'>
type(it): <class 'list_iterator'>
type(first): <class 'list'>
first: ['Jerry', 'Elaine', 'Kramer']
has append on first: True
has __next__ on first: False
\end{verbatim}

The materialized list is a normal list. It is not an iterator, and it does not continue the original traversal state.

The practical distinction is:

\begin{itemize}
    \item Store an iterable when you need a reusable traversal source.
    \item Store an iterator when you intentionally need one moving traversal state.
    \item Convert an iterator to a list or tuple when you need to preserve all remaining values.
\end{itemize}

A final example shows the difference clearly:

\begin{verbatim}
source = range(1, 4)
moving = iter(source)

print("using the source twice")
print(list(source))
print(list(source))

print("using the iterator twice")
print(list(moving))
print(list(moving))
\end{verbatim}

Possible output:

\begin{verbatim}
using the source twice
[1, 2, 3]
[1, 2, 3]
using the iterator twice
[1, 2, 3]
[]
\end{verbatim}

The \texttt{range} object is reusable. The iterator is consumable.

\begin{quote}
An iterable is a source of traversal. An iterator is a traversal in progress. Once that traversal reaches the end, it is exhausted.
\end{quote}